
\documentclass[aspectratio=169]{beamer}
\usetheme{metropolis}           % Use metropolis theme
\usepackage[utf8]{inputenc}
\usepackage{graphicx}
\usepackage{eso-pic}
\usepackage{graphics}
\usepackage{tikz}
\usepackage[export]{adjustbox}
\usepackage{multicol}
\usepackage{listings}
\usepackage{helvet}
\usepackage{booktabs}
\usepackage{threeparttable}
\usepackage{fontspec}
\usepackage{hyperref}
\hypersetup{urlcolor=DarkBlue}

\title{Topic 4 - Track 1 \newline Real Time Data Quality Checks}
\date{\today}
\author{Author of Session here!} % Name of author(s) of session here
\institute{Development Impact Evaluation (DIME) \newline The World Bank }
\setbeamercolor{background canvas}{bg=white}	% Sets background color

% The below command places the World Bank logo and DIME logo to the right corner
\titlegraphic{%
	\begin{picture}(0,0)
	\put(330,-180){\makebox(0,0)[rt]{\includegraphics[width=3cm]{../img/WB_logo}}}
	\end{picture}%
	\begin{picture}(0,0)
	\put(390,-180){\makebox(0,0)[rt]{\includegraphics[width=1.5cm]{../img/i2i}}}
	\end{picture}%
}

%%% Section page with picture of Light bulb
\makeatletter
\defbeamertemplate*{section page}{mytheme}[1][]{
	\centering
	\begin{minipage}{22em}
		\raggedright
		\usebeamercolor[fg]{section title}
		\usebeamerfont{section title}
		\par
		\ifx\insertsubsectionhead\@empty\else%
		\usebeamercolor[fg]{subsection title}%
		\usebeamerfont{subsection title}%
		\fi
		\ifstrempty{#1}{}{%
			\includegraphics[width=100mm, height=60mm]{#1}%
		}
		\insertsectionhead\\[-1ex]
		\insertsubsectionhead
		\usebeamertemplate*{progress bar in section page}
		
	\end{minipage}
	\par
	\vspace{\baselineskip}
}
\makeatother

%%% Define a command to include picture in section, 
%%% make section, and revert to old template
\newcommand{\sectionpic}[2]{
	\setbeamertemplate{section page}[mytheme][#2]
	\section{#1}
	\setbeamertemplate{section page}[mytheme]
}

%%% The command below allows for the text that contains Stata code
\lstset{ %
	backgroundcolor=\color{white},
	basicstyle=\tiny,
	breakatwhitespace=false,
	breaklines=true,
	captionpos=b,
	commentstyle=\color{mygreen},
	escapeinside={\%*}{*)},
	extendedchars=true,
	frame=single,
	numbers=left,
	numbersep=5pt,
	numberstyle=\tiny\color{gray},
	rulecolor=\color{black},
	showspaces=false,
	showstringspaces=false,
	showtabs=false,
	stringstyle=\color{mymauve},
	tabsize=2,
	title=\lstname,
	morekeywords={not,\},\{,preconditions,effects },
	deletekeywords={time}
}

%% The below command creates the ligh bulb logos in the top right corner of the 
\begin{document}
	
	{
		\usebackgroundtemplate{\includegraphics[height=55mm, right]{../img/top_right_corner.pdf}}
		\maketitle
	}

%%%%%%%%%%%%%%%%%%%%%%%%%%%%%%%%%%%%%%%%%%% heading of section 1
\begin{frame}{Objectives}
	\begin{itemize}
		\item Duplicates
			\begin{itemize}
				\item Run a prepared file using ‘ieduplicates’ command and understand the output
			\end{itemize}
		\item Survey Log
			\begin{itemize}
				\item Provide the team with clear progress reports
			\end{itemize}
		\item Back checks 
			\begin{itemize}
				\item Set up bcstats, run it and interpret the results
			\end{itemize}
		\item Other Manual HFC
			\begin{itemize}
				\item Run a prepared HFC file and learn to edit it
			\end{itemize}
	\end{itemize}
\end{frame}

\begin{frame}{ieduplicates}
\begin{itemize}
	\item Exercise 1
	\begin{itemize}
		\item install ieduplicates
		\item learn how to use ieduplicates
		\item practice identifying and correcting duplicates
	\end{itemize}
\end{itemize}
\end{frame}

\begin{frame}{HFC}
\begin{itemize}
	\item Given that you are in track 1, you are usually not expected to know how to write this, so we will focus on how to interpret
	\item We will not cover everything as that is responsibility of person writing
	\begin{itemize}
		\item There are templates that help with this
			\begin{itemize}
				\item IPA
				\item DIME is working
			\end{itemize}
	\end{itemize}
\end{itemize}
\end{frame}


\begin{frame}{3. Survey Log}
\begin{itemize}
	\item Make sure to keep records in the field that is updated very day with how many interviews that were completed
	\item Make sure that the number of observations on the server matches these records.
	\item Purpose: 
		\begin{itemize}
			\item Provide your team a quick overview of progress on the field
			\item Detect enumerators who are slacking
			\item Check balance if this is important for the survey (e.g. by gender)
		\end{itemize}
	\end{itemize}
\end{frame}


\begin{frame}{3. Survey Log}
\begin{itemize}
	\item Exercise 2 – Practice checking survey progress. Assume you’ve verified that the server matches the field log, i.e. no submissions are missing. 
\end{itemize}
\end{frame}


\begin{frame}{4. Back checks}
\begin{itemize}
	\item Purpose
		\begin{itemize}
			\item To monitor the quality of field work: Do some enumerators need extra support from supervisors?
			\item To understand whether your questionnaire accurately captures the key outcomes of your study: Do respondents understand the questions the way we intended to?
		\end{itemize}
	\item Best practices:
		\begin{itemize}
			\item ~ 10\% of surveys, 20\% in the first 2 weeks of field work
			\item Every team and every surveyor must be back checked 
			\item The back check sample must include a proportional number of missing and replacement respondents. 
			\item Selection of households for back checks must be random
		\end{itemize}
\end{itemize}
\end{frame}

\begin{frame}{Selecting Back Check Questions}
\begin{itemize}
	\item Type 1 variables 
		\begin{itemize}
			\item Straightforward questions where we expect very little variation
			\item E.g. education level, marital status, occupation, has children
		\end{itemize}
	\item Type 2 variables
		\begin{itemize}
			\item Questions where we expect capable enumerators to get the true answer
		\end{itemize}
	\item Type 3 variables
		\begin{itemize}
			\item Questions that we expect to be difficult. 
			\item Want to understand if these questions were interpreted in the field as intended		
		\end{itemize}
	\item Identifying Respondent and Interview Information
		\begin{itemize}
			\item Check if we have the right person; if the interview took place \& when
		\end{itemize}
\end{itemize}
\end{frame}

\begin{frame}{bcstats}
\begin{itemize}
	\item ssc install bcstats, replace
	\item Purpose:
	\begin{itemize}
		\item compares back check data and survey data, producing a data set of comparisons
		\item completes enumerator checks and stability checks for variables
	\end{itemize}
\end{itemize}
\end{frame}

\begin{frame}{Using bcstats report to improve data quality}
\begin{itemize}
	\item Type 1
		\begin{itemize}
			\item If >10\% discrepancies  - give surveyor a warning
			\item 20-30\% discrepancies - 2nd backcheck to determine who made errors
			\item If error is by surveyor  - audit 3 additional surveys by that surveyor in same week 
			\item if 20-30\% discrepancies -  put him/her on probation
			\item >40\% discrepancies - 2nd backcheck to determine who made errors and maybe re-survey HH 
			\item If surveyor made errors - re-survey HH \& audit all surveys done by surveyor in this batch
			\item If one more survey has 40\% discrepancies - fire surveyor immediately \& redo all his/her surveys with 20\% or more discrepancies
	\end{itemize}
\end{itemize}
\end{frame}

\begin{frame}{Using bcstats report to improve data quality}
\begin{itemize}
	\item Type 2
		\begin{itemize}
			\item >10\% discrepancies -> Consider re-training
			\item If one surveyor is responsible for more than 30\% of the errors in a single question, follow the steps for Type 1
		\end{itemize}
	\item Type 3
		\begin{itemize}
			\item Discuss >10\% with your survey team and your PIs. 
			\item PIs may decide to edit survey or add additional rounds
		\end{itemize}
\end{itemize}
\end{frame}

\begin{frame}{4. Backchecks}
\begin{itemize}
	\item Exercise 3
		\begin{itemize}
			\item Install bcstats
			\item Learn how to use bcstats
			\item Distinguish between types of backcheck variables 
			\item Explore bcstats output
		\end{itemize}
\end{itemize}
\end{frame}

%%%%%%%%%%%%%%%%%%%%%%%%%%%%%%%%%%%%%%%%%%% Final thougts section
\begin{frame}{Conclusion}


\vspace{20mm}
For more information or further questions please contact:
\newline John Doe (\url{johndoe@worldbank.org}) \newline Mary Doe (\url{marydoe@worldbank.org})

\end{frame}

%%%%%%%%%%%%%%%%%%%%%%%%%%%%%%%%%%%%%%%%%%% The End
\sectionpic{Thank You!}{../img/section_slide}






\end{document} 